\chapter{easy-gl in Depth}
\label{ch:easygl-deep-dive}

This book introduced \texttt{easy-gl} as the OpenGL/OpenGLES wrapper behind CNA's
\texttt{EASYGL} backend. This chapter goes one level deeper --- into its own layering, its
capability-detection mechanism, and a real dependency this book's overview did not have room
to mention.

\section{A layer this book did not cover: meta-gl}

\texttt{easy-gl}'s own README frames it as toolkit-independent, meaning independent of a
windowing toolkit like SDL or Qt --- but it is not independent of a lower C\texttt{++} GL
wrapper underneath it. \texttt{easy-gl} itself is built on a sibling project,
\texttt{meta-gl}, which provides the actual procedural, typed, ownership-free OpenGL/GLES
call surface (typed enum-class wrappers, \texttt{std::span}-based views, lightweight typed
handles, no RAII of its own); \texttt{easy-gl} then layers real object-oriented, RAII,
move-only ownership semantics on top. The dependency is real and direct in the build system
(\texttt{meta-gl} is pulled in as a sibling \texttt{add\_subdirectory} and linked publicly),
and in the source (19 of \texttt{easy-gl}'s own header files include \texttt{meta-gl} headers
directly). The layering is worth remembering precisely because it clarifies what
``toolkit-independent'' actually means here: independent of the host windowing/event toolkit,
not independent of every other C\texttt{++} abstraction below it.

\section{Device: initialization and the GLES-versus-desktop-GL detection mechanism}

\cnaclass{easygl::Device} is the largest single entry point, covering initialization, the
bulk of GL state (clear, viewport/scissor, blend --- including indexed, per-attachment blend
functions for multiple render targets --- depth/stencil, culling, polygon/line mode, sample
coverage), pixel readback (a bounds-checked \texttt{read\_pixels\_robust} alongside the plain
form), the full draw-call family (instanced, indirect, ranged, base-vertex), compute
dispatch, tessellation patch configuration, and debug facilities (labeled debug groups, a
settable debug callback).

\texttt{Device::initialize()}'s own GLES-versus-desktop-GL detection is a genuinely careful
piece of code, not a simple string check: it calls \texttt{glGetString(Version)} and looks
for the substring \texttt{"OpenGL ES"} to decide the API kind, then reads the version number
via the modern \texttt{glGetIntegerv(MAJOR\_VERSION/MINOR\_VERSION)} query --- falling back to
manually parsing the version string on an older context that predates that query, and
specifically accounting for an ES-specific ten-character prefix offset (skipping past the
literal \texttt{"OpenGL ES "} text) that a desktop-GL version string does not have. Extension
enumeration branches the same way: GL~3.0-and-newer contexts use the modern, indexed
\texttt{glGetStringi(Extensions, i)} query; older contexts fall back to splitting the legacy,
single space-delimited \texttt{glGetString(GL\_EXTENSIONS)} string. Once this detection
completes, \texttt{initialize()} hard-requires vertex array objects, shaders, programs,
buffers, and basic rendering support, throwing an \texttt{easygl::Exception} immediately if
any one of them is missing --- refusing to hand back a partially-functional device rather
than letting a caller discover a missing capability later, mid-frame.

\section{Capability gating: a real per-feature table, not a version-number heuristic}

\cnaclass{easygl::Capabilities::detect\_common\_features()} implements a genuine,
per-feature gating table rather than a single ``desktop GL is capability X, GLES is capability
Y'' heuristic, and the two API families frequently diverge in specific, non-obvious ways.
Vertex array objects need GL~3.0 (or the \texttt{GL\_ARB\_vertex\_array\_object} extension) on
desktop, but GLES~3.0 or the differently-named \texttt{GL\_OES\_vertex\_array\_object}
extension on GLES. Framebuffer objects need GL~3.0 or an extension on desktop, but are simply
always present on GLES~2.0 and above, since FBOs are core to ES2 from the start. Tessellation
shaders need GL~4.0 on desktop but only GLES~3.2 --- a case where the GLES requirement is
numerically \emph{lower}. Direct State Access, by contrast, is not modeled for GLES at all,
since it is a desktop-only OpenGL~4.5 capability with no GLES equivalent whatsoever. This
per-feature table --- not a single ``is this GLES'' branch --- is exactly the mechanism
Chapter~\ref{ch:easygl-backend} described CNA's own \texttt{EASYGL} backend
relying on via \texttt{device.supports(...)} / \texttt{device.require(...)}.

\begin{lstlisting}
// Worked example: querying the per-feature table directly, rather than
// branching on Device::capabilities().is_opengles() and guessing.
if (device.supports(easygl::Feature::TessellationShader))
{
    ConfigurePatchPipeline();
}
else
{
    UseFallbackSubdivision();  // e.g. desktop GL < 4.0, or GLES < 3.2
}

// Device::require() throws immediately if the feature is genuinely
// missing, matching Device::initialize()'s own "fail loudly at startup,
// not mid-frame" posture for its hard-required baseline capabilities.
device.require(easygl::Feature::FramebufferObject);
\end{lstlisting}

\section{Real design fixes, and one genuine breaking change worth knowing}

\texttt{easy-gl}'s own planning documents record several real, reasoned design corrections.

\texttt{Program::uniform\_block\_index()} originally returned GL's own raw
\texttt{GL\_INVALID\_INDEX} sentinel on a failed lookup, forcing every caller to know and
compare against a raw GL constant; it now returns \texttt{std::optional<unsigned int>}
instead, so a caller never needs to include a GL header at all just to check for failure.
Several resource classes (\cnaclass{Framebuffer}, \cnaclass{ProgramPipeline},
\cnaclass{TransformFeedback}) were changed from accepting a raw \texttt{unsigned int} GL
handle to accepting a real \cnaclass{Texture} / \cnaclass{Program} reference directly, removing
a whole class of ``passed the wrong handle type by mistake'' bug at the API boundary.

\begin{lstlisting}
// Worked example: both fixes in the same worked snippet. No raw GL
// constant comparison, and no raw GLuint handle to mix up with an
// unrelated resource's handle by mistake.
if (std::optional<unsigned int> blockIndex = shaderProgram.uniform_block_index("Lighting"))
{
    shaderProgram.set_uniform_block_binding(*blockIndex, /* bindingPoint */ 0);
}
// else: no "Lighting" uniform block in this program -- handled without
// ever needing to know or compare against GL_INVALID_INDEX directly.

Framebuffer gbuffer;
gbuffer.attach_texture_2d(FramebufferTarget::Framebuffer, FramebufferAttachment::Depth,
                           TextureTarget::Texture2D, depthTexture, /* level */ 0);
// depthTexture is a real Texture&, not a bare GLuint -- passing a
// Renderbuffer's handle here by mistake is now a compile error, not a
// runtime GL_INVALID_OPERATION discovered later.
\end{lstlisting}

One specific, genuine breaking behavioral change is worth knowing if you are porting code
against an older version of this library: \texttt{Device::clear()} used to call
\texttt{glDisable(GL\_SCISSOR\_TEST)} before every clear, silently overriding whatever scissor
state the caller had already set. This is no longer true --- \texttt{clear()} no longer
touches scissor state at all, and a caller that was relying on the old implicit disable must
now call the scissor-disable method explicitly. The project's own migration notes frame this
as low-impact, since it only matters when scissor testing happens to already be active at the
moment \texttt{clear()} is called --- but it is exactly the kind of silent, implicit
side-effect this book has flagged elsewhere (Chapter~\ref{ch:spritebatch}'s ignored transform
matrix, for instance) as the class of bug most likely to go unnoticed until specifically
tested for.

\subsection{Two small, real RAII additions worth knowing about}

Two small utility classes, added after this chapter's own original pass, are worth a direct
mention because each closes a real, easy-to-get-wrong ergonomic gap rather than adding a new
GL capability. \cnaclass{ScopedDebugGroup} is a two-call RAII pair --- its constructor calls
\texttt{Device::push\_debug\_group}, its destructor unconditionally calls
\texttt{pop\_debug\_group} --- specifically so a labeled debug region in a RenderDoc or NVIDIA
Nsight capture can never be left unbalanced by an early return or an exception unwinding
through the scope:

\begin{lstlisting}
{
    ScopedDebugGroup group(device, "Shadow Pass");
    RenderShadowCasters();
    // pop_debug_group() always runs here, even if RenderShadowCasters()
    // throws or returns early -- no separate try/finally needed.
}
\end{lstlisting}

\cnaclass{UniformCache} addresses the other classic OpenGL ergonomics problem: repeatedly
calling \texttt{glGetUniformLocation} by name is a real, measurable per-draw cost real engines
avoid by caching the resolved integer location once. \texttt{operator[]} looks up a cached
location by name, falling through to a real \texttt{Program::uniform\_location} call and
memoizing the result only on a cache miss --- and, worth knowing explicitly since it is not
stated in the header's own brief comment, a \emph{failed} lookup (a typo'd or since-optimized-out
uniform name, returning $-1$) is memoized exactly the same way a successful one is. Calling
\texttt{operator[]} again with the same wrong name will not re-attempt the GL query and
magically start succeeding once the shader changes --- only an explicit \texttt{invalidate()}
call (needed after any program relink, since old locations are not guaranteed valid against a
newly linked program) clears the cache and allows a fresh lookup:

\begin{lstlisting}
UniformCache uniforms(shaderProgram);
shaderProgram.set_uniform(uniforms["uTime"], elapsedSeconds);
shaderProgram.set_uniform(uniforms["uLightDir"], lightDir.x, lightDir.y, lightDir.z);
// ... later, only after shaderProgram is actually relinked ...
uniforms.invalidate();  // required -- old cached locations are not safe to reuse otherwise
\end{lstlisting}

\section{Testing without a GPU}

\texttt{easy-gl}'s own test suite is three small, hand-rolled (not GoogleTest-based)
\texttt{assert()}-driven executables, and the mechanism making them useful without any real
GPU or display is worth describing directly: a fake \texttt{GLGetProcAddressFn} loader
returns hand-written function pointers only for the handful of GL entry points
\texttt{Device::initialize()} actually calls, returning literal hardcoded strings (for
instance, \texttt{"OpenGL ES 3.0"}) and literal hardcoded enum values in place of real driver
responses. This lets the entire \cnaclass{ContextInfo} / \cnaclass{Capabilities} detection
pipeline run and be asserted against genuinely, entirely off-GPU --- every other GL function
pointer the loader might be asked for resolves to a non-null but non-callable sentinel, so a
test can exercise detection logic without ever risking an actual GPU call.

\section{Two known, honestly stated limitations}

Two gaps are worth knowing about directly rather than assuming closed.

First, \texttt{Config::enable\_debug\_logging} currently does nothing at all --- no
\texttt{log\_callback} field yet exists on \texttt{Config} to actually receive the output,
despite the flag existing to request it. Second, \cnaclass{Query::result\_u64()} is
implemented using the 32-bit \texttt{glGetQueryObjectuiv} internally, because the underlying
\texttt{meta-gl} layer does not expose the real 64-bit query variant --- meaning it is
suitable for frame counts, but not for a high-resolution GPU timer query, despite its
64-bit-suggesting name. Both are stated plainly in the project's own tracking rather than
silently left for a reader to discover.
